\subsubsection{Codage des couples tritiques et somme avec retenue}
\label{mg04:codificacion}

La lecture ternaire des émissions APP--TRIT--TPK a pour espace ambiant
six coordonnées sur \(\mathbb F_3\). Leur regroupement en trois couples
ordonne cet espace ambiant au moyen de trois chiffres nonadiques. Le regroupement
conserve la position et l'ordre des trits ; la somme qu'il représente
doit aussi transporter sa retenue.

Soit \(T=\{0,1,2\}\) l'ensemble des représentants ternaires. Pour une classe \(x\in\mathbb Z/9\mathbb Z\), nous écrivons \(\widetilde x\in\{0,\ldots,8\}\) pour son représentant canonique. Nous définissons
\begin{equation}
 \beta:T^2\longrightarrow\mathbb Z/9\mathbb Z,
 \qquad \beta(a,b)=[a+3b]_9.
 \label{mg04:beta}
\end{equation}
Ici, \([s]_m\) désigne, lorsqu'il est utilisé comme chiffre, le représentant
dans \(\{0,\ldots,m-1\}\) de la classe de \(s\).

\begin{proposition}[Inverse et transport de la somme]
L'application \(\beta\) est une bijection, d'inverse
\begin{equation}
 \beta^{-1}(x)=
 \left([\widetilde x]_3,
       \left\lfloor\frac{\widetilde x}{3}\right\rfloor\right).
 \label{mg04:inversa}
\end{equation}
La loi sur \(T^2\) transportée depuis l'addition nonadique est
\begin{equation}
 (a,b)\oplus(c,d)=
 \left([a+c]_3,
 \left[b+d+\left\lfloor\frac{a+c}{3}\right\rfloor\right]_3\right).
 \label{mg04:suma}
\end{equation}
Avec cette loi, \(\beta\) est un isomorphisme de groupes additifs.
L'élément neutre est \((0,0)\) et l'inverse de \((a,b)\) est
\begin{equation}
 \left([-a]_3,
 \left[-b+\left\lfloor\frac{-a}{3}\right\rfloor\right]_3\right).
 \label{mg04:opuesto}
\end{equation}
\end{proposition}

\begin{proof}
La division euclidienne de chaque \(0\leq\widetilde x\leq8\) par
trois donne un unique reste \(a\in T\) et un unique quotient \(b\in T\),
avec \(\widetilde x=a+3b\). Cela prouve la bijection et
\eqref{mg04:inversa}.

Pour l'addition, écrivons \(a+c=3k+r\), avec \(r=[a+c]_3\) et \(k=\lfloor(a+c)/3\rfloor\). Alors
\[
 (a+3b)+(c+3d)=r+3(b+d+k).
\]
La réduction de \(b+d+k\) modulo trois ne change le membre droit que d'un multiple de neuf. Ainsi, \(\beta((a,b)\oplus(c,d))=\beta(a,b)+\beta(c,d)\). La bijectivité transporte l'associativité, la commutativité et l'élément neutre. Enfin, la division \(-a=3\lfloor-a/3\rfloor+[-a]_3\) appliquée à \(-a-3b\) donne \eqref{mg04:opuesto}.
\end{proof}

La retenue est déjà visible dans
\[
 (2,0)\oplus(1,0)=(0,1),\qquad
 \beta(2,0)+\beta(1,0)=[3]_9.
\]
L'addition vectorielle de \(\mathbb F_3^2\) produirait \((0,0)\)
dans ce même exemple. De même, \((1,0)\) est d'ordre neuf pour
\(\oplus\), tandis que son ordre pour l'addition vectorielle est trois.
Ces opérations spécifient les deux structures algébriques qui
coexistent sur le même ensemble de neuf représentants.

\begin{corollary}[Coordinatisation cubique de l'espace ambiant ternaire et fibres de projection]
L'application coordonnée par coordonnée
\begin{equation}
 \begin{split}
 \beta^{\times3}:T^6&\longrightarrow(\mathbb Z/9\mathbb Z)^3,\\
 (t_1,t_2,t_3,t_4,t_5,t_6)&\longmapsto
 \bigl(\beta(t_1,t_2),\beta(t_3,t_4),\beta(t_5,t_6)\bigr)
 \end{split}
 \label{mg04:cubo}
\end{equation}
est une bijection d'ensembles de \(729=3^6=9^3\) éléments. La
projection \(\operatorname{pr}_{12}(X,Y,Z)=(X,Y)\) sur les
\(81\) positions nonadiques possède exactement neuf éléments
par fibre.
\end{corollary}

\begin{proof}
Le produit des trois inverses de \eqref{mg04:inversa} inverse \eqref{mg04:cubo}. \(X,Y\) étant fixées, la troisième coordonnée parcourt les neuf classes de \(\mathbb Z/9\mathbb Z\) ; d'où \(\operatorname{pr}_{12}^{-1}(X,Y)=
\{(X,Y,Z):Z\in\mathbb Z/9\mathbb Z\}\) et le cardinal affirmé. Conserver \(Z\) avec \(X,Y\) permet d'appliquer l'inverse complète et de récupérer les six trits ordonnés.
\end{proof}

Le domaine de ce dénombrement est l'espace ambiant complet de \(729\)
mots. Le catalogue de \(468\) émissions décimales et son image
admissible de \(243\) mots ternaires conservent leurs propres applications
émettrices et leurs restrictions. La restriction de \(\beta^{\times3}\)
à cette image demeure injective, tandis que les cardinaux de
ses fibres de projection sont déterminés par l'intersection avec la
fibre ambiante. L'affirmation uniforme de neuf éléments appartient
au cube complet.

La coordinatisation bidimensionnelle obtenue restitue une représentation étiquetée
du support APP déjà construit. Les feuillets de somme et de produit sont des fonctions supplémentaires sur
ce support, avec les représentants et quotients déclarés dans le noyau.
La troisième coordonnée conserve la donnée omise par la projection
faciale ; les routes et la mémoire temporelle demeurent dans l'état TPK
dont provient la lecture. Le codage fini
explicite ainsi l'information qu'il transporte et la distingue de l'historique
complet auquel il appartient.
